\uuid{Dckd}
\exo7id{3668}
\titre{exo7 3668}
\auteur{quercia}
\organisation{exo7}
\datecreate{2010-03-11}
\isIndication{false}
\isCorrection{false}
\chapitre{Espace euclidien, espace normé}
\sousChapitre{Produit scalaire, norme}
\cpge{mpsi-mp2i.espaces-prehilbertiens.a, pcsi.espaces-prehilbertiens.a, ptsi.geometrie-plan.b, mpsi-mp2i.espaces-prehilbertiens.b, pcsi.espaces-prehilbertiens.b, ptsi.geometrie-espace.b, mpsi-mp2i.espaces-prehilbertiens.e, pcsi.espaces-prehilbertiens.e, mp-mpi.topologie-espaces-normes.a, pc.espaces-vectoriels-normes.a, psi.espaces-vectoriels-normes.a}
\module{Algèbre}
\niveau{L2}
\difficulte{}

\contenu{
\question{
Soit $E$ un espace vectoriel euclidien et $p \in \mathcal{L}(E)$ une projection.
Montrer que :

$p$ est une projection orthogonale 
$\Leftrightarrow \forall\ \vec x,\vec y \in E,\ (\vec x\mid p(\vec y)) = (p(\vec x)\mid \vec y)$\par
$\Leftrightarrow \forall\ \vec x \in E,\ \|p(\vec x)\| \le \|\vec x\,\|$.

(Pour la deuxième caractérisation, considérer $\vec x \in (\mathrm{Ker} p)^\perp$ et
faire un dessin)
}
}
